%%% Vyplňte prosím základní údaje o závěrečné práci (odstraňte \xxx{...}).
%%% Automaticky se pak vloží na všechna místa, kde jsou potřeba.

% Druh práce:
%	"bc" pro bakalářskou
%	"mgr" pro diplomovou
%	"phd" pro disertační
%	"rig" pro rigorozní
\def\ThesisType{bc}

% Název práce v jazyce práce (přesně podle zadání)
\def\ThesisTitle{{Herní prostředí pro hru inspirovanou hrou Logix a tvorba ukázkových botů, s důrazem na trénování ohodnocovací funkce pro MCTS}}

% Název práce v angličtině
\def\ThesisTitleEN{{Game environment for a game inspired by the game Logix and development of sample bots with emphasis on training of evaluation function for MCTS}}

% Jméno autora (vy)
\def\ThesisAuthor{{František Procházka}}

% Rok odevzdání
\def\YearSubmitted{{2026}}

% Název katedry nebo ústavu, kde byla práce oficiálně zadána
% (dle Organizační struktury MFF UK:
% https://www.mff.cuni.cz/cs/fakulta/organizacni-struktura,
% případně plný název pracoviště mimo MFF)
\def\Department{{Katedra teoretické informatiky a matematické logiky}}
\def\DepartmentEN{{Department of Theoretical Computer Science and Mathematical Logic}}

% Jedná se o katedru (department) nebo o ústav (institute)?
\def\DeptType{{Katedra}}
\def\DeptTypeEN{{Department}}

% Vedoucí práce: Jméno a příjmení s~tituly
\def\Supervisor{{Mgr. Vladan Majerech, Dr.}}

% Pracoviště vedoucího (opět dle Organizační struktury MFF)
\def\SupervisorsDepartment{{Katedra teoretické informatiky a matematické logiky}}
\def\SupervisorsDepartmentEN{{Department of Theoretical Computer Science and Mathematical Logic}}

% Studijní program (kromě rigorozních prací)
\def\StudyProgramme{{Umělá inteligence}}

% Nepovinné poděkování (vedoucímu práce, konzultantovi, tomu, kdo
% vám nosil pizzu a vařil čaj apod.)
\def\Dedication{%
{Rád bych poděkoval svému vedoucímu práce Mgr. Vladanu Majerechovi, Dr., za odborné vedení, cenné připomínky a vstřícnost při směřování této práce.}
}

% Abstrakt (doporučený rozsah cca 80-200 slov; nejedná se o zadání práce)
\def\Abstract{%
{Tato práce se zabývá návrhem a implementací herního prostředí pro abstraktní deskovou hru inspirovanou hrou Logix a tvorbou 
ukázkových botů pro toto prostředí. Součástí práce je formální popis pravidel hry, reprezentace herního stavu, návrh množiny legálních 
akcí a implementace prostředí umožňujícího simulaci partií. Dále práce popisuje návrh botů založených na jednoduchých heuristikách, 
metodě Monte Carlo Tree Search a algoritmu Alfa-beta ořezávání. Zvláštní důraz je kladen na přípravu trénovací pipeline pro ohodnocovací funkci, 
která může být využita při řízení prohledávání v MCTS. Výsledkem práce je funkční prostředí vhodné pro experimentování s různými strategiemi 
a základ pro další vývoj silnějších hráčů založených na metodách strojového učení.}
}

% Anglická verze abstraktu
\def\AbstractEN{%
{This thesis focuses on the design and implementation of a game environment for an abstract board game inspired by Logix and on the
 development of sample bots for this environment. The work includes a formal description of the game rules, representation of the game state,
 design of the set of legal actions, and implementation of an environment capable of simulating complete games. The thesis also describes sample
  bots based on simple heuristics, Monte Carlo Tree Search and Alpha-Beta pruning. Special emphasis is placed on the preparation of a training 
  pipeline for an evaluation function that can be used to guide the MCTS search. The result is a functional environment suitable for experimenting
   with different strategies and a foundation for further development of stronger players based on machine learning methods.}
}

% 3 až 5 klíčových slov (doporučeno) oddělených \sep
% Hodí se pro nalezení práce podle tématu.
\def\ThesisKeywords{%
{desková hra\sep herní prostředí\sep Monte Carlo Tree Search\sep ohodnocovací funkce\sep strojové učení}
}

\def\ThesisKeywordsEN{
{board game\sep game environment\sep Monte Carlo Tree Search\sep evaluation function\sep machine learning}
}

% Pokud některá z položek metadat obsahuje TeXové řídící sekvence, je potřeba
% dodat i verzi v obyčejném textu, která se objeví v metadatech formátu XMP
% zabudovaných do výstupního souboru PDF. Pokud si nejste jistí, prohlédněte si
% vygenerovaný soubor thesis.xmpdata.
\def\ThesisAuthorXMP{\ThesisAuthor}
\def\ThesisTitleXMP{\ThesisTitle}
\def\ThesisKeywordsXMP{\ThesisKeywords}
\def\AbstractXMP{\Abstract}

% Máte-li dlouhý abstrakt a nechceme se mu vejít na informační stranu,
% můžete použít toto nastavení ke zmenšení písma informační strany.
% (Uvažte nicméně zkrácení abstraktu, to je často lepší.)
\def\InfoPageFont{}
%\def\InfoPageFont{\small}  % odkomentujte pro zmenšení písma
